% Terminal apparatus of the concise study: the scope note and the References
% for the sources this companion actually uses. The revision timestamp and the
% rights colophon follow from the shared generation record.

\section{Appendix: Scope and Qualifications}

{\footnotesize
\textbf{Edition, formulary and occurrence.} This concise study treats the Mass
\latin{Dominica vigesima prima post Pentecosten}, II classis, of the \work{Missale
Romanum}, editio typica, Typis Polyglottis Vaticanis 1962, pp.~413--415,
nos.~1699--1708, in the universal 1962 calendar and no particular calendar. Its
occurrence on 18~October 2026, its precedence over the feast of St Luke, the
commemoration of that feast in low and conventual Masses and its omission in a
sung Mass that is not conventual, the oration \latin{pro Fidei propagatione} added
in every Mass on the penultimate Sunday of October, the colour, the Credo and the
Trinity Preface rest on the Missal's own calendar and general rubrics (RG 16,
108, 111, 127~b; RGMR 434, 437~a, 444--446, 453, 475~a, 494~b). Whether a local
Ordinary has appointed another day for the missions, and any particular calendar,
are local facts not applied here. The added orations are not part of the
formulary and are not studied here.

\textbf{Relation to the full study.} Every claim here is drawn from the full
study of this same formulary, the Full PDF subtitled \textit{The war God bounds,
the debt forgiven from the heart, and the Church's exodus: a study of the proper
in three readings}, which prints the appointed Latin with its public-domain
English, treats each element in its own setting, and develops the three readings
at length, each with its own four senses; nothing evidence-dependent is added
here. The four rows on the first page orient a reader and do not replace those
readings' own senses. Each Father and later writer is cited only for what he said
of his own passage, and each liturgical commentator only for what he says of the
elements of this formulary he names. Bl.\ Ildefonso Schuster and the continuation of \work{The Liturgical Year} comment
on this formulary as the 1962 Missal gives it.

\textbf{Text, English and rights.} Every Latin form of the ten elements was read
for the study on the page images of the CMAA facsimile of the typical edition and
collated with the Pustet \work{Missale Romanum} of Ratisbon, 1862, in which each
element's wording already stands; the collation is in
\texttt{propers/verified.md}. No witness read outside liturgical books carries the
Introit's antiphon in the Missal's form, the Missal's opening of the Communion or
its added \latin{Domine Deus meus}, or the wording of the Offertory's narrative; no
chant manuscript, Old Latin or Roman psalter was consulted. Psalms are cited in the
Missal's Vulgate numbering, with the Hebrew numbering added in the map. The
scriptural English is the Douay--Rheims in Challoner's revision at the canonical
verses, so that it does not render the Missal's Latin where a chant departs from
its verse; the English of the Introit's first sentence, the Offertory, the
Communion and the three orations is the anonymous translation printed by Cummiskey
in Philadelphia in 1861. Neither is an approved liturgical translation of the 1962
Missal. The English of the Fathers, Doctors and commentators is that of the
published translations named below; English beside the Latin of St Jerome, St
Hilary, St Thomas or Cassiodorus is a gloss of that Latin. The Latin of the
formulary rests on its public-domain antecedent, the Pustet Missal (17 U.S.C.
103(b)); the Bibles, the 1861 English, Migne's texts and the other editions and
translations quoted are in the public domain in the United States. The \work{New
American Bible} introductions are summarised, not reproduced.

\textbf{Reception and its limits.} The witnesses named here are those the study's
sweep read at the loci below. St Hilary's commentary on Ps~118 and St Jerome's on
Matthew and Ephesians were read on the page images of Migne's \work{Patrologia
Latina}, and the exposition of the psalms printed under St Athanasius's name on
the page images of Migne's \work{Patrologia Graeca}, whose attribution of that
work is the printing's; the other witnesses were read in text layers,
transcriptions or web texts. St Thomas names Origen as holding the view that
forgiven sins return, which he denies. Among the witnesses not reached are St
Ambrose on Ps~118, Greek commentators on Pss~89 and 118 other than that
exposition, Theophylact and Cornelius a Lapide on Ephesians, Hugh of Saint-Cher,
and St Anthony of Padua's Sunday sermons.

\textbf{Chronology and records.} The rows on the second page come from the
Scripture chronology corpus through the record generated for this formulary: its
default profile's answer for every element and, for the Gospel and the Epistle, a
declared comparison under its critical profile. The record holds no occasion for
any of the psalms, no attribution for Ps~89 or for Job, and no date for the
telling of the parable; the Davidic attribution of Ps~118 rests on Haydock's note,
and the reign shown is the reference point of that attribution. The search, its
loci, the disagreements and the negative results are in
\texttt{research/scope.md}, the reasoning behind the three readings in
\texttt{research/interpretations.md}, and this companion's production record in
\texttt{research/production-review.md}.
\par}

\phantomsection
\section*{References}
\runninghead{References}

{\scriptsize
\subsection*{Liturgical books and Bibles}
\begin{itemize}\setlength{\itemsep}{0.15em}\setlength{\parskip}{0pt}
\item \work{Missale Romanum ex decreto Sacrosancti Concilii Tridentini restitutum, Summorum Pontificum cura recognitum}, editio typica (Typis Polyglottis Vaticanis, 1962), CMAA facsimile, \url{https://media.churchmusicassociation.org/pdf/missale62.pdf}: \work{Rubricae generales} 16, 108, 111, 127; \work{Rubricae generales Missalis} 434, 437, 444--446, 453, 475, 494; \work{Ordo Missae} no.~1130; pp.~413--415 (nos.~1699--1708).
\item \work{Missale Romanum} (Ratisbon: Pustet, 1862), pp.~354--355, \latin{Dominica XXI. post Pentecosten}.
\item \work{The Roman Missal, Translated into the English Language for the Use of the Laity} (Philadelphia: Eugene Cummiskey, 1861), pp.~450--452 (XXI Sunday after Pentecost: Introit, Collect, Offertory, Secret, Communion, Postcommunion).
\item The Holy Bible, Douay--Rheims version, revised by Bishop Richard Challoner (Project Gutenberg text): 3~Kgs 6:1; Esth 1:2; 13:8--18; Job 1:1, 12, 21; 9:4; Ps 89:1--2; 113:1--3; 118:1, 81--86; Lam 4:21; Mt 17:23; 18:1, 23--35; 20:2; Lk 10:19; Rom 11:29; 13:14; 16:20; 1~Cor 10:13; Eph 6:10--17, 20; Apoc 6:10--11. \work{Biblia Sacra Vulgatae editionis} (Clementine Vulgate), eBible.org text: Esth 13:9--11; Job 1:1, 12; Ps 89:1--2; 113:1; 118:1, 81--86; Mt 18:23--35; Eph 6:10--17.
\end{itemize}

\subsection*{Fathers, Doctors and saints}
\begin{itemize}\setlength{\itemsep}{0.15em}\setlength{\parskip}{0pt}
\item St Augustine, \work{Enarrationes in Psalmos}, trans. in Nicene and Post-Nicene Fathers, 1st ser., vol.~8 (CCEL text): on Ps~89:1--3; 113:3--4; 118, sermon~1.3 and \latin{Caph} (vv.~84, 86).
\item St Augustine, \work{Sermons} 58 and 83 (Benedictine numbering; Sermons 8 and 33 in the translation) and \work{De sermone Domini in monte} I.22.77, trans. in Nicene and Post-Nicene Fathers, 1st ser., vol.~6 (CCEL text).
\item St Gregory the Great, \work{Morals on the Book of Job} (\work{Moralia in Iob}), Library of the Fathers (Oxford): Preface 7, 16; II.19; IX.5.
\item St Jerome, \work{Commentarii in Matthaeum} III, on Mt 18:23--35, and \work{Commentarii in Epistolam ad Ephesios} III, on Eph 6:11--12, in Migne, PL 26 (Paris, 1845), cols.~132--133 and 543--546.
\item St Hilary of Poitiers, \work{Tractatus super Psalmos}, in Ps.~118, \latin{Caph}~5, in Migne, PL 9 (Paris, 1844), cols.~573--574.
\item St John Chrysostom, \work{Homilies on Matthew} 61, Nicene and Post-Nicene Fathers, 1st ser., vol.~10 (CCEL text); \work{Homilies on Ephesians} 22--24, Nicene and Post-Nicene Fathers, 1st ser., vol.~13 (CCEL text).
\item St Thomas Aquinas, \work{Summa theologiae} I q.~103 a.~7, trans. Fathers of the English Dominican Province; \work{Super Evangelium S.~Matthaei lectura} c.~18, in \work{Opera}, editio altera Veneta, vol.~3 (Venice, 1745), pp.~242--244; \work{In Epistolam ad Ephesios} c.~6, lect.~3--4, in \work{In omnes D.~Pauli Apostoli epistolas commentaria}, vol.~2 (Li\`ege: H.~Dessain, 1857), pp.~356--359; \work{Expositio super Iob ad litteram} c.~1, Corpus Thomisticum web text.
\item St Rabanus Maurus, \work{Commentaria in Matthaeum}, on Mt 18:23--35, in Migne, PL 107 (text layer).
\item St Robert Bellarmine, \work{A Commentary on the Book of Psalms}, trans. J.~O'Sullivan, eCatholic2000 transcription: on Pss~89:1; 118:84.
\item St Alphonsus Liguori, \work{Sermons for All the Sundays in the Year}, trans. N.~Callan (Dublin: James Duffy, 1882), Sermon L, pp.~380--387.
\item Cassiodorus, \work{Expositio Psalmorum}, on Pss~113:1; 118:1 and 84, Corpus Corporum web text.
\item \work{Expositiones in Psalmos}, printed under the name of St Athanasius, in Migne, PG 27 (Paris): on Ps~118:84--86, col.~495.
\end{itemize}

\subsection*{Liturgical commentators}
\begin{itemize}\setlength{\itemsep}{0.15em}\setlength{\parskip}{0pt}
\item Bl.\ Ildefonso Schuster, \work{The Sacramentary (Liber Sacramentorum)}, vol.~3 (London: Burns Oates \& Washbourne, 1927), pp.~178--181.
\item The continuation of Prosper Gu\'eranger, \work{The Liturgical Year} (Dom Lucien Fromage), \work{The Time after Pentecost}, series vol.~11, trans. Laurence Shepherd (Dublin: James Duffy and Sons, 1883), pp.~456--470.
\item Rupert of Deutz, \work{De divinis officiis} XII.21 and \work{De omnipotentia Dei} 9, in Migne, PL 170 (text layer).
\item Honorius Augustodunensis, \work{Gemma animae} IV.93, in Migne, PL 172 (Paris, 1854).
\item William Durandus, \work{Rationale divinorum officiorum} VI.138--139, vol.~2 (Lyon, 1612; text layer).
\end{itemize}

\subsection*{Chronology sources}
\begin{itemize}\setlength{\itemsep}{0.15em}\setlength{\parskip}{0pt}
\item \work{The Catholic Encyclopedia} (New York, 1907--1912): ``Biblical Chronology'' (vol.~3, 1908); ``David'' (vol.~4, 1908); ``Esther'' and ``Epistle to the Ephesians'' (vol.~5, 1909); ``Gospel of Saint Matthew'' (vol.~10, 1911); ``St.\ Paul'' (vol.~11, 1911); ``The New Testament'' (vol.~14, 1912).
\item G.~L. Haydock, \work{The Douay--Rheims Bible with Haydock's Commentary}, note on Ps~118:1.
\item United States Conference of Catholic Bishops, \work{New American Bible, Revised Edition}, introductions to Job, the Psalms, Matthew and Ephesians, on their dating; summarised, not reproduced.
\end{itemize}
\par}
